\subsection{Finite embedding problems and the obstruction left by ordinary inertia}

\begin{result}{E.140}{Proved}{The finite odd-Tate embedding theorem and its exact scope}
Let \(\ell\) be odd and let a finite embedding problem over
\(\mathbb Q\) have abelian kernel
\[
                    A=(\mathbb Z/\ell^n(w))^b
\]
with the indicated actual Galois action. If it is soluble
at every completion, it has a global weak solution.
Given a global weak solution, any prescribed local solutions
at a finite set \(S\) can be attained while retaining that
solution at a disjoint finite set \(T\). A proper solution
can also be obtained with these prescriptions.

Here solutions are considered up to conjugation by the
kernel. The theorem fixes the finite lower homomorphism.
It supplies neither a specified ramification set outside
\(S\cup T\), nor integral Tate-module lifts, nor ordinary
point-section behaviour at the newly ramified primes.
\end{result}
\begin{proof}
For one cyclic summand, finite first-cohomology restriction
to any finite set is surjective by
\citet[Theorem 2.1(e), author p.4; Corollary 2.3,
author p.6]{CohNeftin2013}. The latter statement retains
a second finite set, and explicitly assumes an existing
global solution. Applying the assertion to each cyclic
summand proves the required surjectivity for \(A\).
Differences of solutions of the fixed embedding problem
are \(A\)-cocycles, as in that paper's setup,
author pp.3--4. Thus the first-cohomology assertion gives
the stated local modification.

We supply the separate weak-solvability argument.
The dual \(A^D=\operatorname{Hom}(A,\mu_{\ell^n})\)
has cyclic Galois action, because
\((\mathbb Z/\ell^n)^\times\) is cyclic.
Every nonzero class in \(H^1(\mathbb Q,A^D)\)
is detected at an unramified finite prime outside any
given finite set. Indeed, over the field splitting
the module it is a homomorphism. If this homomorphism
is nonzero, Chebotarev detects it at a prime splitting
in that field. If it is zero, the class is inflated
from the finite cyclic action group; a prime whose
Frobenius generates that group detects it.
Inflation from that cyclic decomposition quotient is
injective, so detection here is detection in local
cohomology.

Let an absolute \(H^2(\mathbb Q,A)\)-class be locally
zero everywhere. Represent it over \(G_{\mathbb Q,U}\)
for a finite \(U\) containing all coefficient and
ramification primes. The group \(H^1(G_{\mathbb Q,U},A^D)\)
is finite. Choose finitely many extra primes detecting
every nonzero element of its local kernel at \(U\).
A class in the corresponding local kernel for the
enlarged set is locally zero at the new primes,
so has zero inertia there and descends to \(G_{\mathbb Q,U}\).
The enlarged first local kernel is therefore zero.
Finite Poitou--Tate duality identifies its dual with
the second local kernel
\citep[I.4.10]{CohMilneADT}. Our second-cohomology
class becomes zero over this enlarged set, hence was
zero absolutely. This proves the needed injectivity
\[
 H^2(\mathbb Q,A)\longrightarrow
                         \prod_v H^2(\mathbb Q_v,A).
\]
Apply it to the actual embedding obstruction. No
projectivity of the absolute Galois group is assumed.

For properness, select \(b\) distinct auxiliary primes
splitting completely in the fixed lower extension
and in \(\mathbb Q(\mu_{\ell^n})\), outside \(S\cup T\).
At the \(i\)-th such prime prescribe a kernel-valued
local solution whose tame inertia is the \(i\)-th
basis vector of \(A\) and whose Frobenius is one.
The tame relation holds because \(q\equiv1\pmod{\ell^n}\).
The same finite first-cohomology theorem attains these
prescriptions together with the original ones.
Its image contains \(A\) and maps onto the lower
group, so is the whole extension group.
The argument permits further ramification and makes
no claim that these auxiliary local solutions are
ordinary point sections.
\end{proof}

\begin{result}{E.141}{Proved}{An explicit finite correction which has no integral lift}
Let \(\phi_7\) be the cubic character of
\(\mathbb Q(\zeta_7)^+/\mathbb Q\), normalized by
\([2]\mapsto1\) in
\((\mathbb Z/7)^\times/\{\pm1\}\simeq\mathbb F_3\).
Regard it as a class in
\(H^1(\mathbb Q,\mathbb F_3(4))\).
Its local Frobenius values at two and five are
\[
                              (1,1).
 \tag{E.141.1}
\]
No integral class in \(H^1(\mathbb Q,\mathbb Z_3(4))\)
has these two reductions. In particular \(\phi_7\)
has a nonzero integral coefficient Bockstein,
already visible locally at seven.
\end{result}
\begin{proof}
The fourth finite Tate twist is trivial modulo three.
Both two and five represent \([2]\) in the displayed
cubic cyclotomic quotient, since \(5=-2\bmod7\).
This proves (E.141.1).
E.67 gives
\[
 H^1(\mathbb Q,\mathbb Z_3(4))=\mathbb F_3\eta_4,
                  \qquad \eta_4=(\chi_3^4-1)/3.
 \]
Its local reductions at two and five are
\[
 \left((2^4-1)/3,\ (5^4-1)/3\right)
                         =(2,1)\quad\text{in }\mathbb F_3^2.
 \tag{E.141.2}
\]
The vector \((1,1)\) is not in this one-dimensional
image. The coefficient sequence therefore gives
a nonzero Bockstein
\(\delta\phi_7\in H^2(\mathbb Q,\mathbb Z_3(4))[3]\).

There is also a direct local check. The character is
tamely ramified at seven, whereas all integral
weight-four one-cocycles there have zero inertia.
The integral local second-cohomology group is
\[
 H^2(\mathbb Q_7,\mathbb Z_3(4))
                  =\mathbb Z_3/(7^3-1)
                  \simeq\mathbb Z/9.
 \]
The ramified finite character maps to its nonzero
subgroup of order three. With a choice of tame
generator it is a unit multiple of
\((7^3-1)/3=114\) modulo \(342\).
Thus this failure of integral lifting is not
inferred merely from a global dimension count.
\end{proof}

\begin{result}{E.142}{Proved}{Which auxiliary primes can change finite and integral central parameters}
Let \(S\) contain three and infinity. For even \(w\)
and \(A=\mathbb Z/3^n(w)\), allowing any additional
set \(T\) of primes \(q\equiv-1\pmod3\) changes neither
finite first nor finite second cohomology:
\[
 H^i(G_{\mathbb Q,S},A)
       \xrightarrow{\ \sim\ }
 H^i(G_{\mathbb Q,S\cup T},A),
                         \qquad i=1,2.
 \tag{E.142.1}
\]
For integral coefficients and \(w\ge2\), first
cohomology is unchanged by allowing \emph{any}
new primes different from three:
\[
 H^1(G_{\mathbb Q,\{3,\infty\}},\mathbb Z_3(w))
       \xrightarrow{\ \sim\ }
 H^1(G_{\mathbb Q,S\cup T},\mathbb Z_3(w))
       \xrightarrow{\ \sim\ }
 H^1(\mathbb Q,\mathbb Z_3(w)).
 \tag{E.142.2}
\]
The second assertion does not extend to integral
second cohomology. In particular auxiliary primes
congruent to one modulo three can supply finite
corrections such as E.141, without supplying any
new integral central parameter.
\end{result}
\begin{proof}
The two residue terms of finite localization at
\(q\ne3\) are
\(H^0(\mathbb F_q,A(-1))\) and
\(H^1(\mathbb F_q,A(-1))\).
Frobenius acts on \(A(-1)\) by \(q^{w-1}\).
For the primes in (E.142.1),
\(q^{w-1}-1\) is a three-adic unit, so both terms
vanish. Finite purity and the \(G_S\) comparison
used in A.22 and E.66 prove (E.142.1) for finite
\(T\). For arbitrary \(T\), every finite continuous
cochain factors through an extension with finite
ramification. Passing through the filtered union
of finite subsets proves the same assertion.

For (E.142.2), the action on the integral module
is unramified at \(q\ne3\). A one-cocycle on
inertia is a homomorphism, and its tame value
must satisfy
\[
                    (q^{w-1}-1)z(T_q)=0 .
\]
This nonzero integer acts injectively on
\(\mathbb Z_3\), even when it is divisible by
three. Wild inertia and the other tame factors
map trivially as well. Hence the cocycle has zero
inertia at every such prime and descends to
\(G_{\mathbb Q,\{3,\infty\}}\).
Coboundaries descend too. This proves the
integral assertion directly, without interchanging
an infinite set of ramification places with a
coefficient inverse limit.
\end{proof}

\begin{result}{E.143}{Proved}{Ordinary inertia forbids the usual new central ramification at fixed exponent}
Let \(Q\) be a finite geometric pro-three quotient
of exponent \(3^a\), whose abelianization is
\((\mathbb Z/3^a)^2\) with the marked peripheral
generators. At \(q\ne3\), an ordinary local point
section whose abelian inertia is zero has trivial
geometric inertia in \(Q\).

Consequently a new nontrivial central inertia
introduced by a finite embedding solution, while
its abelian inertia remains zero, cannot be
absorbed by choosing nonzero ordinary integer
valuations divisible by \(3^a\).
This obstruction occurs already in the marked
Heisenberg quotient \(H_3\).
\end{result}
\begin{proof}
An ordinary point with both valuations zero has
good reduction away from the three punctures;
its prime-to-\(q\) path representation is
unramified
\citep[Definition 2.2 and Proposition 2.3, p.200]
{AuditWojtkowiak2007}.
Otherwise its tame inertia is conjugate to a
power of \(x\), \(y\), or the infinity peripheral
generator, with exponent its nonzero ordinary
valuation, as in A.31.
Zero abelian inertia forces that exponent to
be divisible by \(3^a\). Every such power is
the identity in the exponent-\(3^a\) group \(Q\).
Conjugacy does not change this conclusion.

For a literal example, take the actual point
section at \(5/16\) in \(H_3\), and multiply it
by the central character \(\phi_7\) of E.141.
The actual action on the center is \(\chi_3^2\),
which is trivial modulo three, so this is a
valid global finite section with the same
abelian reduction. The point \(5/16\) has good
reduction at seven: its two coordinates reduce
to six and two. The new section therefore has
nontrivial central inertia at seven and zero
abelian inertia there. Since \(H_3\) has
exponent three, the first assertion excludes
every ordinary local point at seven, including
points with nonzero valuations divisible by
three. This is a bad solution of a soluble
finite embedding problem; it does not assert
that all its solutions are bad.
\end{proof}

\begin{result}{E.144}{Proved}{A necessary frozen-support gate for the changing-pair embedding route}
Let \(F=\Pi^{(3)}\) and use its lower exponent-three
central series
\[
 P_1=F,\qquad P_{n+1}=\overline{P_n^3[P_n,F]} .
 \]
For \(n\ge2\) put
\[
 Q_{n-1}=F/P_n,\qquad Q_n=F/P_{n+1},\qquad
 W_n=P_n/P_{n+1},\qquad
 B_n=\langle x^{3^{n-1}},y^{3^{n-1}}\rangle
                         \ \subset W_n .
 \tag{E.144.1}
\]
Then \(B_n\simeq\mathbb F_3(1)^2\).
The actual action on \(W_n\) factors through
\(\chi_3\bmod3\), and hence
\[
 W_n/B_n=C_n^+\oplus C_n^-,
 \qquad C_n^+\simeq\mathbb F_3^{\,a_n},\quad
 C_n^-\simeq\mathbb F_3(1)^{\,d_n}.
 \tag{E.144.2}
\]

Retain the actual action on \(Q_n\) in a finite
arithmetic quotient \(D\), unramified outside
three. Fix a lower homomorphism
\[
 \rho:G_{\mathbb Q,S}\longrightarrow Q_{n-1}\rtimes D
 \]
with its prescribed arithmetic projection.
Suppose it matches specified ordinary point
sections at every old place in \(S\).
Any extension of \(\rho\) to \(Q_n\rtimes D\)
which retains these old local point conditions,
and whose local class at each newly ramified
prime is the image of an ordinary point section,
must induce a solution
\[
 \bar\rho:G_{\mathbb Q,S}\longrightarrow
                         (Q_n/B_n)\rtimes D
 \tag{E.144.3}
\]
which retains all these old local point
conditions. This necessity is unchanged if
arbitrarily many new primes are allowed,
including primes congruent to one modulo three.

For \(S=\{3,\infty,q_1,\ldots,q_t\}\), with all
\(q_i\equiv-1\pmod3\), the quotient embedding in
(E.144.3) has an unramified global weak solution.
Its old local matching is an exact finite
first-cohomology localization problem for
\(C_n^+\oplus C_n^-\). More explicitly,
\[
 \begin{split}
 H^1(G_{\mathbb Q,S},\mathbb F_3)
       &=\mathbb F_3\cdot\bar\eta_4,\\
 H^1(G_{\mathbb Q,S},\mathbb F_3(1))
       &=\langle\kappa(3),\kappa(q_1),\ldots,
                                      \kappa(q_t)\rangle .
 \end{split}
 \tag{E.144.4}
\]
For each choice of old local aligning frames,
comparison with one global weak solution gives
a finite linear system in these stated parameters.
Allowing all finite aligning frames gives the
precise quotient matching condition. No
sufficiency for a full ordinary \(Q_n\)-extension,
or for an infinite tower, is asserted.
\end{result}
\begin{proof}
Each \(W_n\) is a finite elementary abelian
three-group, central in the geometric quotient
\(Q_n\). The two displayed peripheral powers
map to a basis of the kernel in abelianizations
\[
 (\mathbb Z/3^n)^2\longrightarrow
                         (\mathbb Z/3^{n-1})^2 .
 \]
They are therefore independent and span the
claimed two-dimensional \(B_n\).
The actual formula for the action on \(x,y\)
shows it is the cyclotomic action on these
powers: geometric conjugation disappears in
the central layer. This proves the assertion
about \(B_n\).

An automorphism trivial on \(F/P_2\) is trivial
on every \(P_n/P_{n+1}\).
This follows inductively from the definition:
replacing an element of \(P_n\) by its product
with an element of \(P_{n+1}\) changes its cube
only in \(P_{n+2}\), and replacing either entry
in a generating commutator changes it only
in \(P_{n+2}\). Here one uses
\([P_i,P_j]\subseteq P_{i+j}\), proved by the
same commutator induction.
The graded pieces are generated by these
cubes and commutators.
Thus the Galois action factors through its
scalar action \(\chi_3\bmod3\) on the Frattini
quotient. Complex conjugation acts by an
involution; its plus and minus eigenspaces
split over \(\mathbb F_3\).
On a commutator expression followed by cube
operations its sign is the parity of the
number of original generators in the
commutator. This proves (E.144.2) without
assuming a diagonal action on the integral
nilpotent quotient.

Let \(q\notin S\) be newly ramified.
Its fixed lower image has trivial inertia.
The lower abelian quotient forces the
ordinary valuation to be a multiple of
\(3^{n-1}\). In \(Q_n\) its point inertia
is consequently in the line generated by
\(x^{3^{n-1}}\), by \(y^{3^{n-1}}\), or by
the corresponding infinity power.
The latter also belongs to \(B_n\):
\[
 (xy)^{3^{n-1}}
       =x^{3^{n-1}}y^{3^{n-1}}
                                      \pmod{P_{n+1}}.
 \tag{E.144.5}
\]
For completeness, the cube map on
\(P_i/P_{i+1}\) is additive into
\(P_{i+1}/P_{i+2}\) at the odd prime three.
For \(i=1\), the extra terms in \((uv)^3\)
are commutator cubes and commutators of
length three, all in \(P_3\).
For \(i\ge2\), the commutator terms already
belong to \(P_{i+2}\).
Iterating this fact proves (E.144.5).
Geometric conjugation acts trivially on
\(W_n\). Hence every newly allowed ordinary
inertia is killed on quotienting by \(B_n\).
The quotient homomorphism is unramified
outside the old \(S\), and its old local
point images are retained. This proves
the necessary assertion including its
Frobenius conditions, not just an inertia
condition.

We justify the asserted finite description.
The embedding obstruction for (E.144.3)
belongs to
\(H^2(G_{\mathbb Q,S},C_n^+\oplus C_n^-)\).
Its local restrictions at \(S\) vanish
because the prescribed actual points lift.
The plus group has zero second cohomology
by E.66 and E.142.
For the minus group, Kummer theory and
\(\operatorname{Pic}\mathbb Z[S^{-1}]=0\)
identify second cohomology with
\(\operatorname{Br}(\mathbb Z[S^{-1}])[3]\).
The Brauer local invariant sequence is
injective on this group. Thus its locally
zero obstruction also vanishes. This
produces a global weak solution over the
old ramification set.

The plus first-cohomology equality follows
from E.142 and E.66. The minus equality
is the \(S\)-unit Kummer sequence; the
listed classes form a basis.
The fibres over a fixed literal lower
homomorphism are torsors for this first
cohomology, modulo conjugation by the
kernel. Compare one weak solution with
the point lifts at the old places.
Their differences must belong to the
image of the explicit localization map
in (E.144.4).
Finite projected stabilizers are not
discarded: use each possible aligning
element of the finite lower geometric
group, and test the resulting affine
linear system. Different lifts of an
aligner change only a kernel coboundary.
This is a finite collection of finite
linear tests.

For example, at \(q_i\ne3\) the minus
local class is its valuation, and the
coefficient of \(\kappa(q_i)\) is thereby
fixed. At three the residual unit
coordinate must equal the sum of those
fixed coefficients times
\(\kappa(q_i)|_{\mathbb Q_3}\); the remaining
\(\kappa(3)\) coefficient controls the
valuation coordinate.
For each plus coordinate there is only
the single cyclotomic cubic parameter,
with its actual local unit factors.
These old constraints cannot be altered
by adding ordinary inertia in \(B_n\)
at new primes. Passing this quotient
test still leaves the \(B_n\)-extension
and all ordinary constraints at its
new primes to be constructed.
\end{proof}

\begin{result}{E.145}{Proved}{The free three-adic Galois group and the relations created by frozen support}
Put \(K=\mathbb Q(\zeta_3)\). The maximal pro-three
quotient of \(G_{K,\{3\}}\) is free pro-three of
rank two.
If \(t\) rational primes \(q_i\) with
\[
 q_i\equiv-1\pmod3,\qquad v_3(q_i^2-1)=1
 \]
are added to the ramification set, its maximal
pro-three quotient \(P_S\) instead satisfies
\[
 \dim_{\mathbb F_3}H^1(P_S,\mathbb F_3)=t+2,
 \qquad
 \dim_{\mathbb F_3}H^2(P_S,\mathbb F_3)=t .
 \tag{E.145.1}
\]
In particular it is not free when any such
ordinary support prime has been frozen.
The new second-cohomology classes are detected
by the explicit cups
\(\kappa(\zeta_3)\cup\kappa(q_i)\).
Thus projectivity at the coefficient prime
alone cannot be applied to the Galois group
which retains this old support.
\end{result}
\begin{proof}
Shapiro's lemma for the quadratic extension
\(K/\mathbb Q\), together with its two
characters over \(\mathbb F_3\), gives
\[
 H^i(G_{K,S_K},\mathbb F_3)
   =H^i(G_{\mathbb Q,S},\mathbb F_3)
      \oplus H^i(G_{\mathbb Q,S},\mathbb F_3(1)).
 \tag{E.145.2}
\]
For \(S=\{3,\infty\}\), E.66 makes second
cohomology zero, and E.67's finite calculation
makes first cohomology two-dimensional.
For \(t\) added primes, E.142 leaves the
plus groups equal to dimensions one and zero.
The \(S\)-unit Kummer sequence gives minus
first-cohomology dimension \(t+1\).
The Brauer invariant sequence gives minus
second-cohomology dimension \(t\): it is
the kernel of the sum map from \(t+1\)
copies of \(\mathbb F_3\) at the finite
places of \(S\). This proves those
dimensions for the full \(G_{K,S_K}\).

For any profinite group \(G\), its maximal
pro-three quotient \(P\) has the same
first cohomology with \(\mathbb F_3\),
and inflation in degree two is injective.
To see the latter, let \(N\) be its kernel.
A \(P\)-invariant homomorphism
\(N\to\mathbb F_3\) would define a quotient
of \(G\) which is an extension of a
pro-three group by a three-group, hence
pro-three. Maximality of \(P\) forces that
homomorphism to be zero.
The five-term sequence proves the injection.

When \(t=0\), second cohomology of \(P\)
therefore vanishes. Here is a direct
freeness argument. Take a minimal
surjection from a free pro-three group
on two generators to \(P\), with kernel
\(R\). The five-term sequence and
vanishing second cohomology imply
\(\operatorname{Hom}(R,\mathbb F_3)^P=0\).
If \(R\ne1\), some finite three-group
quotient has nonzero image \(R_0\).
The finite normal subgroup \(R_0\)
has a nonzero quotient
 \(R_0/(R_0^3[R_0,G_0])\):
apply the elementary local-ring Nakayama
lemma to its Frattini quotient as a
module for the finite three-group \(G_0\).
It gives a forbidden invariant character
of \(R\). Thus \(R=1\).

For \(t>0\), each rational \(q_i\) is
inert in \(K\), with residue field
\(\mathbb F_{q_i^2}\).
The classes of \(\zeta_3\) and \(q_i\)
are \(S_K\)-unit Kummer classes and
factor through \(P_S\).
Their cup has nonzero tame residue at
the prime over \(q_i\):
the residue of \(\zeta_3\) is not a
cube, since \(v_3(q_i^2-1)=1\).
At each other newly allowed prime both
entries are units, so the residue is
zero. These \(t\) cups are therefore
linearly independent.
They span the full second-cohomology
group just computed, and each is in
the image of inflation from \(P_S\).
Inflation, already injective, is thus
an isomorphism in this case.
This proves (E.145.1) and its claimed
nonprojective obstruction.
\end{proof}
